% Re-scoped abstract (2026-09-18): leads with the phenomenon, presents the theory as a solvable model, and does not claim that the
% random-feature mechanism explains the trained MLPs. Matches the abstract registered on OpenReview. abstract.tex is left untouched.
\begin{abstract}
Latent diffusion models operate in a latent space whose dimension $\dlat$ typically exceeds the intrinsic dimension $\dint$ of the data. We show that this excess delays memorization. For MLP score networks of fixed width trained on synthetic mixtures with known intrinsic dimension, the fraction of memorized samples after 5M steps falls from $30\%$ at $\dlat=5$ to near zero at $\dlat=40$, and the divergence of the learned score from the population score begins progressively later as the latent widens. The decline is not an artifact of nearest-neighbor tests losing sensitivity in high dimension: it survives a dimension-controlled test restricted to the true signal subspace. On CelebA and CIFAR-10 VAE latents, once the VAE is wide enough to preserve sample identity, the time to memorize grows two- to four-fold with latent width, while FID follows a separate quality tradeoff. To understand the effect we evaluate the random-feature spectral theory of \citet{bonnaire2025} at the two-atom covariance of data embedded in a larger latent space. The feature-correlation spectrum splits into signal, noise-dimension, sample-specific and null blocks; under gradient flow the delay between generalization and memorization is the ratio of two spectral edges, not a count of intervening modes; and this ratio grows with $\dlat$ because extra coordinates lower the pre-activation variance and shrink the sample-specific feature mass, which a norm-fixed control confirms. We then test which parts of this mechanism carry over to trained networks, and use the latent spectrum to build a one-parameter clock for the memorized fraction during training.
\end{abstract}
